% =============================================================================
% CHAPTER 2 — LITERATURE SURVEY
% =============================================================================
\chapter{Literature Survey}
\label{ch:lit}

\section{Introduction}

A structured review of published literature has been conducted to situate the StockSense AI
system within the broader research landscape of inventory management, demand forecasting,
artificial intelligence applications in retail, and workflow automation. The review spans
journal articles, conference papers, and technical reports drawn from IEEE Xplore, arXiv, and
other indexed databases. The goal of the survey is to identify what has been achieved in related
domains, what gaps remain unaddressed, and how the proposed system is positioned to fill those
gaps.

\section{Review of Existing Systems}

\subsection{Intelligent Inventory Systems for Small and Micro Enterprises}

Wang et al.\ (2024) \cite{wang2024} present an intelligent inventory system for small and micro
enterprise warehousing that combines RFID-based tracking with a rule-based alerting engine
deployed on a cloud platform. The system demonstrates measurable reductions in manual counting
errors. However, the reliance on RFID hardware infrastructure creates a significant financial
barrier, as the cost of tagging individual stock-keeping units across an entire kirana store is
impractical. The system also does not address supplier evaluation or AI-based reorder
recommendations.

\subsection{Demand Forecasting with Machine Learning for Multi-Channel Retail}

Kheawpeam and Sinthupinyo (2023) \cite{kheawpeam2023} propose a demand forecasting model using
machine learning techniques including LSTM and gradient boosting applied to a multi-channel
retail dataset. The study shows that LSTM models outperform traditional ARIMA baselines in
handling non-linear seasonal demand patterns. While technically rigorous, the model requires
substantial historical data (at minimum twelve months) and dedicated model retraining
infrastructure. These prerequisites are unrealistic for kirana stores that may lack even basic
transaction records.

\subsection{Retail Demand Forecasting with External Influences}

Soni and Vaghela (2024) \cite{soni2024} conduct a comparative study of machine learning models
for retail demand forecasting that explicitly incorporates external variables such as weather,
festivals, and promotional events. The study demonstrates improvements in forecast accuracy
when contextual features are included. The work is relevant to StockSense AI's future scope
of seasonal demand multipliers; however, the current implementation of this paper is limited
to academic model comparison without a deployable system for resource-constrained environments.

\subsection{LLM Applications in Business Decision-Making}

Al Khaldy and Gheraibia (2024) \cite{alkhaldy2024} evaluate the capabilities and limitations of
large language models in business decision support contexts, including procurement and inventory.
The study finds that LLMs can reduce cognitive bias in reorder decisions by providing consistent,
data-driven reasoning. However, the paper does not address the integration of LLMs into
operational real-time workflows or the cost implications of using LLM APIs at the scale of
small businesses.

\subsection{Mediating Cognitive Biases Using LLMs}

Friedman, Krug, and Mountford (2025) \cite{friedman2025} extend this line of inquiry by
investigating how LLMs can be used to explicitly counteract documented cognitive biases in
business purchasing decisions, such as anchoring bias with preferred suppliers. The work
validates the suitability of LLM-generated reasoning as a decision-support layer, which directly
supports the design choice in WF-05 of StockSense AI to use Groq Llama~3.1~8B for reorder
suggestion generation.

\subsection{Workflow Automation Efficiency Using n8n}

Amir and Atif (2025) \cite{amir2025} present the first academic evaluation of n8n as a workflow
automation platform in a small-scale business context. The study finds that n8n enables
non-programmers to construct business automation pipelines with significantly lower development
effort compared to custom-coded solutions, while maintaining comparable execution performance.
This finding directly validates the architectural decision in StockSense AI to use n8n as the
backend orchestration engine.

\subsection{IoT-Based Smart Inventory Management}

Manoharan et al.\ (2024) \cite{manoharan2024} propose an IoT-based smart inventory system that
uses weight sensors and RFID to trigger automated reorder alerts. The system demonstrates
high accuracy in detecting low-stock conditions but requires continuous power supply and
hardware maintenance, making it operationally impractical for small unassisted retail stores.
The work underscores the hardware dependency problem that software-only solutions like
StockSense AI are designed to circumvent.

\subsection{Smart Retail Using ML for Demand Forecasting and Inventory Optimisation}

Poongothai et al.\ (2024) \cite{poongothai2024} describe a smart retail framework combining
convolutional neural networks for product recognition with LSTM-based demand forecasting. The
framework achieves high prediction accuracy but presupposes camera infrastructure, high-bandwidth
connectivity, and trained operations staff. The work illustrates the performance ceiling
achievable with complex ML pipelines while confirming the need for simpler, deployable
alternatives for the unorganised retail segment.

\section{Comparative Analysis}

Table~\ref{tab:litcomp} presents a structured comparison of the reviewed works against StockSense
AI across key evaluation dimensions relevant to small Indian retail deployment.

\begin{table}[H]
\centering
\caption{Comparative Analysis of Existing Solutions vs.\ StockSense AI}
\label{tab:litcomp}
{\singlespacing
\renewcommand{\arraystretch}{1.3}
\begin{tabularx}{\textwidth}{@{}l X X X X@{}}
\toprule
\textbf{Feature} & \textbf{StockSense AI} & \textbf{Vyapar} & \textbf{Zoho} & \textbf{Khatabook} \\
\midrule
AI Reorder & \textbf{Yes} & No & Limited & No \\
Stockout Prediction & \textbf{Automated} & Threshold alerts only & Basic & No \\
Supplier Scoring & \textbf{WSM-based} & No & Basic & No \\
WhatsApp Alerts & \textbf{Automated} & Transactional only & Transactional only & No \\
Monthly Cost & \textbf{Rs.\ 0--500} & Rs.\ 0--599 & Rs.\ 2999+ & Rs.\ 0 \\
Setup Complexity & \textbf{Low} & Medium & High & Low \\
Real-time Dashboard & \textbf{Yes} & Limited & Yes & Basic \\
\bottomrule
\end{tabularx}
}
\end{table}

\section{Research Gap and Motivation}

The literature review reveals a clear and consistent pattern: existing inventory management
research falls into two distinct clusters. The first cluster comprises academically rigorous
machine learning and IoT-based systems that achieve high predictive performance but require
substantial hardware investment, skilled personnel, and complex infrastructure that is entirely
inaccessible to small Indian retailers. The second cluster comprises commercial low-cost tools
such as Khatabook or basic POS applications that are financially accessible but provide only
passive transaction recording without any predictive, prescriptive, or AI-driven intelligence.

The specific combination of capabilities present in StockSense AI — automated stockout prediction,
LLM-generated reorder reasoning, multi-criteria supplier scoring, AI vision for data entry,
real-time WhatsApp alerts, and full pipeline orchestration — does not appear as a unified system
in any reviewed work. The research gap is therefore not in any single capability but in their
synthesis into a coherent, low-cost, deployable package for the kirana store context.

Furthermore, the reviewed literature contains no published work on using n8n as a production
backend for an AI-augmented inventory system, and no work on combining Groq LLM APIs with
Supabase real-time subscriptions for live inventory dashboards. StockSense AI addresses both
of these novel integration challenges and demonstrates their practical viability.

The motivation for this project is therefore grounded in both a social need — making inventory
intelligence accessible to India's twelve million kirana store owners — and a technical gap in
the published literature that the proposed system is uniquely positioned to address.
